\documentclass[10pt, oneside]{article}
\usepackage{geometry}

%\usepackage[T1]{fontenc}
\usepackage{geometry}
\usepackage{changepage} 
\usepackage[pdftex]{graphicx}
\usepackage{multicol}
\usepackage{ gensymb }
\usepackage[english,spanish]{babel}
\usepackage[utf8]{inputenc}
\usepackage{enumitem} 
\setlist{leftmargin=3mm}

\usepackage{dcolumn}

\makeatletter 

\newlength\tdima
\newcommand\tabfill[1]{%
      \setlength\tdima{\linewidth}%
      \addtolength\tdima{\@totalleftmargin}%
      \addtolength\tdima{-\dimen\@curtab}%
      \parbox[t]{\tdima}{#1\ifhmode\strut\fi}}

\newcommand\mytabs{\hspace*{1cm}\=\hspace{1cm}\=\hspace{1cm}\=\hspace{1cm}\=\hspace{1cm}\=\hspace{1cm}\=\hspace{1cm}\=\hspace{1cm}\=\hspace{1cm}\=\hspace{1cm}}
\newenvironment{mysec}[1][\mytabs]
  {\begin{tabbing}#1\kill\ignorespaces}
  {\end{tabbing}}
 
\makeatother
  
\usepackage{titlesec}

\newcommand{\lyxrightaddress}[1]{
\par {\raggedleft \begin{tabular}{l}\ignorespaces
#1
\end{tabular}
\vspace{1.4em}
\par}
}

\geometry{letterpaper}
\geometry{verbose,tmargin=1.3cm,bmargin=1.3cm,lmargin=1.2cm,rmargin=1.2cm}
% Activate to surpress page number on first page:
\thispagestyle{empty}

\usepackage{soul}
\usepackage{times}

\usepackage{ifthen,xcolor}
\newlength{\tabcont}
\newcommand{\tab}[1]{%
\settowidth{\tabcont}{#1}%
\ifthenelse{\lengthtest{\tabcont < .25\linewidth}}%
{\makebox[.25\linewidth][l]{#1}\ignorespaces}%
}%

\usepackage{hyperref}
\hypersetup{colorlinks=true, urlcolor=black}
 
\urlstyle{same}

\usepackage{titlesec}
\titleformat{\subsection}
  {\large\scshape\bfseries}
  {\thesubsection}{1em}{}

\setlength{\columnsep}{0pt}
\hypersetup{
     colorlinks   = true,
     urlcolor    = blue
}
\begin{document}

%%%%%%%%%%%%%%%%%%%%%%%%%%%%%%%%%%%%%%%%%%%%%%%%%%%%%%%%%%%%%%%%%%%%%%%%%%%%%%%%%%%%%%%%%%%%%%%%%%%%%%%%%%%%

\pagenumbering{gobble}

\begin{center}

\begin{picture}(1000,1)
    \put(0,-20){\includegraphics[width=\textwidth]{img/formalheader.png}}
\end{picture}

\vspace{10mm}
{ University of California, Berkeley \hfill Department of Agricultural and Resource Economics} \\\hrule  \vspace{10mm}
{\Large \textbf{Lucy Hackett}} \\
\end{center} 

%%% CONTACT INFORMATION
\begin{minipage}[t]{0.1\linewidth}
\textbf{Contact \\ Information}
\end{minipage}\hspace{0.05\linewidth}
\begin{minipage}[t]{0.8\linewidth}
Giannini 241, UC Berkeley \\
\href{mailto:lucy_hackett@berkeley.edu}{lucy\_hackett@berkeley.edu} \\
\href{https://lghackett.github.io}{https://lghackett.github.io} \\
+1-(510)-326-3451
\end{minipage}\vspace{5mm}

%%% DOCTORAL STUDIES
\begin{minipage}[t]{0.1\linewidth}
\textbf{Doctoral Studies}
\end{minipage}\hspace{0.05\linewidth}
\begin{minipage}[t]{0.8\linewidth}
University of California, Berkeley\\
PhD, Agricultural and Resource Economics, Expected completion May 2026 \\
\textsc{Dissertation:}  ``Essays in Environmental and Urban Economics'' \\ 
~\\ 
\textsc{Primary Fields:} Environment \\
\textsc{Secondary Fields:} Development, Urban/Spatial Economics \\
\vspace{0.3cm} 

\begin{minipage}[t]{0.5\linewidth}
\begin{itemize}[noitemsep,nolistsep]
\item[] \underline{Associate Professor Marco Gonzalez-Navarro} \\
\href{mailto:marcog@berkeley.edu}{marcog@berkeley.edu} \\ 
+1 (510) 390-4720 \\
Department of Agricultural
\item[] \hspace{4mm} \& Resource Economics \\
\item[] \underline{Professor Lucas Davis} \\
\href{mailto:lwdavis@berkeley.edu}{lwdavis@berkeley.edu} \\
%+1 (999) Phone \\
Haas School of Business 
\end{itemize}
\end{minipage}
\begin{minipage}[t]{0.4\linewidth}
\begin{itemize}[noitemsep,nolistsep]
    \item[] \underline{Assistant Professor Kirill Borusyak} \\
\href{k.borusyak@berkeley.edu}{k.borusyak@berkeley.edu} \\ 
%+1 (999) Phone \\
Department of Agricultural
\item[] \hspace{4mm} \& Resource Economics\\\\ 
\item[] \underline{Professor Joseph Shapiro} \\
\href{mailto:joseph.shapiro@berkeley.edu}{joseph.shapiro@berkeley.edu} \\ 
%+1 (999) Phone \\
Department of Agricultural
\item[] \hspace{4mm} \& Resource Economics
\end{itemize}
\end{minipage}
\end{minipage}\vspace{5mm}

%%% PLACEMENT OFFICERS
\begin{minipage}[t]{0.1\linewidth}
\textbf{Placement\\ Officers}
\end{minipage}\hspace{0.04\linewidth}
\begin{minipage}[t]{0.9\linewidth}~\vspace{-9mm}
\begin{multicols}{3}
\begin{itemize}[noitemsep,nolistsep]
\item[] \underline{Professor Sofia Villas-Boas} \\
\href{mailto:sberto@berkeley.edu}{sberto@berkeley.edu} \\
+1 (510) 409-4341 \\
\item[] \underline{Professor Marco Gonzalez-Navarro} \\
\href{mailto:marcog@berkeley.edu}{marcog@berkeley.edu} \\
+1 (510) 390-4720 \\
\item[] \underline{Diana Lazo} \\
\href{mailto:lazo@berkeley.edu}{lazo@berkeley.edu} \\
+1 (510) 642-3345 
\end{itemize}
\end{multicols}
\end{minipage}\vspace{5mm}

%%% PRIOR EDUCATION
\begin{minipage}[t]{0.1\linewidth}
\textbf{Prior \\ Education}
\end{minipage}\hspace{0.05\linewidth}
\begin{minipage}[t]{0.8\linewidth}
\begin{mysec} 
\textbf{University of Oregon} \>\>\>\>\> B.Sc. in Economics and Spanish Literature \` 2015 
\end{mysec} 
\begin{mysec} 
    \textbf{CIDE} \>\>\>\>\> Master's Degree in Economics \` 2018
    \end{mysec} 
\end{minipage}\vspace{5mm}

%%% TEACHING
\begin{minipage}[t]{0.1\linewidth}
\textbf{Teaching}
\end{minipage}\hspace{0.05\linewidth}
\begin{minipage}[t]{0.8\linewidth}
\begin{mysec} 
    \textbf{UC Berkeley} \>\>\>Instructor, \emph{Microeconomics}  \` 2022, 2023 \\
    \>\>\> Master's in Development Practice
    \end{mysec}
\begin{mysec} 
\textbf{UC Berkeley} \>\>\>Teaching Assistant, \emph{Graduate Econometrics}  \` 2021, 2022, 2023 \\
\>\>\> PhD Agricultural and Resource Economics
\end{mysec}
\begin{mysec} 
\textbf{UC Berkeley} \>\>\>Instructor, \emph{Training for New Graduate Student Instructors}  \` 2021, 2022, 2023 \\
\>\>\> UC Berkeley 
\end{mysec}
\begin{mysec} 
    \textbf{CIDE} \>\>\>Instructor, \emph{Graduate Macroeconomics}  \` 2020 \\
    \>\>\> Master's in Economics 
\end{mysec}
\begin{mysec} 
    \textbf{CIDE} \>\>\>Instructor, \emph{Macroeconomics}  \` 2019 \\
    \>\>\> Economics (undergraduate)
\end{mysec}
\begin{mysec} 
    \textbf{CIDE} \>\>\>Teaching Assistant, \emph{Macroeconomics}  \` 2017, 2018, 2019 \\
    \>\>\> Economics (Graduate and undergraduate)
\end{mysec}
\end{minipage}\vspace{5mm}

%%% LANGUAGES
\begin{minipage}[t]{0.1\linewidth}
\textbf{Languages}
\end{minipage}\hspace{0.05\linewidth}
\begin{minipage}[t]{0.8\linewidth}
English (native), Spanish (fluent)
\end{minipage}\vspace{5mm}

%%% GRANTS, FELLOWSHIPS, AWARDS
\begin{minipage}[t]{0.1\linewidth}
\textbf{Grants, Fellowships, and Awards}
\end{minipage}\hspace{0.05\linewidth}
\begin{minipage}[t]{0.8\linewidth}
\begin{mysec} 
	2025 \>\>\tabfill{UC Dissertation Fellowship (\$37,000)} \\
    2025 \>\>\tabfill{CEGA Development Economics Challenge Grant (\$5,000)} \\
    2024 \>\>\tabfill{Clausen Center for International Trade Research Grant (\$5,000)} \\
    2024 \>\>\tabfill{Fisher Center for Real Estate Research Grant  (\$15,000)} \\
	2023 \>\>\tabfill{Fisher Center for Real Estate Research Grant  (\$15,000)} \\
    2023 \>\>\tabfill{Giannini Foundation Minigrant  (\$35,000)} \\
    2022 \>\>\tabfill{ARE travel grant  (\$2,000)} \\
    2021 \>\>\tabfill{Tinker Field Research Grant (\$2,000)}\\ 
    2021 \>\>\tabfill{Outstanding Graduate Student Instructor Award}\\ 
    2018 \>\>\tabfill{Best applied thesis, Class of 2018, CIDE}\\ 
    2015 \>\>\tabfill{Phi Beta Kappa}\\ 
    2015 \>\>\tabfill{Oregon 6}
\end{mysec} 
\end{minipage}\vspace{5mm}

%%% RESEARCH PAPERS
\begin{minipage}[t]{0.1\linewidth}
\textbf{Research \\ Papers}
\end{minipage}\hspace{0.05\linewidth}
\begin{minipage}[t]{0.8\linewidth}
\textbf{``Land Subsidence: Environmental risk in housing markets in Mexico City''\\ 
(\textsc{Job Market Paper})}\\ 
 with Carolina Rodriguez-Zamora \\

We study the costs of and the housing market response to subsidence, the sinking of land areas due to groundwater over-extraction, in Mexico City. 
%Subsidence is a prevalent and worsening phenomenon worldwide. 
%We estimate the costs of both plot-level subsidence and uneven neighborhood sinking through the housing market, quantifying the damages to private housing units and public infrastructure. 
We propose an equilibrium model of the housing market that features housing re-development in the face of an evolving environmental hazard that has both realized and expected future impacts to home quality. 
%Our model highlights that while realizations of subsidence attract development by lowering the opportunity cost of re-building units, expected future sinking deters development if predictable risk is capitalized into home values. 
Guided by model-derived estimating equations for key parameters of the model, we  exploit quasi-random variation in sinking intensity to estimate the impact of both realized and future subsidence on home values. We find that realized subsidence imposes substantial costs, which are driven by physical degradation to the structure, increased maintenance investment, and impacts to public infrastructure. However, prices are unresponsive to measures of expected future sinking, and novel survey evidence on residents' beliefs and information about sinking suggest that information frictions affect the ability of homebuyers to capitalize predictable future risk. Units that have experienced more sinking are more likely to be redeveloped, as these have lower opportunity cost of being re-built. Evaluating welfare using our parameter estimates implies that subsidence costs Mexico City \$34.2 billion USD in economic costs, 12\% of which are due to information frictions that inefficiently increase the housing stock in risky areas. Our findings show that groundwater depletion imposes a costly externality on the built environment, and that frictions affecting the capitalization of environmental hazards in the housing market exacerbate these costs by putting more value in harm's way. \\

\textbf{``Estimating the Gains from Water Trade: A Systematic
Evaluation of Modeling Considerations''} with Nell Green Nylen, Ellen Bruno, Andrew Ayers, Michael Kiparsky, Josué Medellín-Azuara, and Sarah Null \\ 
\emph{Conditionally accepted at the Review of Environmental Economics and Policy. Draft available upon request.} \\

The gains from water trading can vary significantly depending on local conditions as well as the specifics of market design and implementation. However, models of water trading necessarily rely on assumptions that simplify the social, institutional, and environmental landscape within which a water market operates. We systematically evaluate peer-reviewed papers that estimate the gains from water trading to assess how models of water markets  take this local context into account. Our results demonstrate that whether and how models incorporate key considerations varies widely, with implications for the accuracy of results. We find that estimates of the economic impacts of water trading in the published literature are more likely to consider distributional effects and incorporate features of the legal and regulatory environment than to account for third-party impacts, transaction costs, the consequences of trading for the economy at large, or the administrative costs associated with setting up and operating a market. Understanding what features a model takes into account is important for interpreting its policy implications. Researchers modeling the gains from trade could better support local decision makers by explicitly articulating their models’ capabilities and limitations.
\end{minipage}\vspace{5mm}


%%% RESEARCH IN PROGRESS
\begin{minipage}[t]{0.1\linewidth}
\textbf{Research in Progress}
\end{minipage}\hspace{0.05\linewidth}
\begin{minipage}[t]{0.8\linewidth}
\textbf{``Differential subsidence, damages and fragility: Evidence from a systematic analysis in Mexico City''} with Enrique Fernández-Torres\\ 

\textbf{``Sprawl and the costs and benefits of promoting compact cities: Evidence from Mexico''} with Rodolfo Oviedo Moguel \\ 

\textbf{``Rural road upgrading and transportation in Rwanda''} with Marco Gonzalez-Navarro, Florence Kondylis, John Loeser, and Maria Ruth Jones\\ 

\textbf{``Protect the forest, protect cities? Land use change externalities and flooding in Brazil''} with Alice Schmitz\\ 

\textbf{``Differential subsidence, damages, and fragility: Evidence from a systematic analysis of structural vulnerability in Mexico City''} with Enrique Fernández-Torres
\end{minipage}

%%% TALKS
\begin{minipage}[t]{0.1\linewidth}
\textbf{Talks}
\end{minipage}\hspace{0.05\linewidth}
\begin{minipage}[t]{0.8\linewidth}
\begin{mysec} 
    2025 \>\> \tabfill{NEUDC (\textit{upcoming})} \\
    2025 \>\> \tabfill{Cities \& Development Workshop, Harvard (\textit{upcoming})} \\
    2025 \>\> \tabfill{The Occasional Workshop of Environmental and Resource Economics (\textit{upcoming})} \\
    2025 \>\> \tabfill{North American Meeting of the Urban Economics Association} \\
    2025 \>\> \tabfill{Environment Week at the London School of Economics} \\
    2025 \>\> \tabfill{AERE Summer meeting} \\
	2024 \>\> \tabfill{LACEA Urban Workshop} \\
    2024 \>\> \tabfill{\textit{SobreMéxico} Conference,  Universidad Iberoamericana}\\ 
    2024 \>\> \tabfill{UC Berkeley Development workshop} \\
    2023, 2024 \>\> \tabfill{UC Berkeley Trade workshop} \\
\end{mysec} 
\end{minipage}\vspace{5mm}

%%% REFEREEING
% \begin{minipage}[t]{0.1\linewidth}
% \textbf{Refereeing}
% \end{minipage}\hspace{0.05\linewidth}
% \begin{minipage}[t]{0.8\linewidth}
% \emph{Journal of Excellent Papers}
% \end{minipage}\vspace{5mm}

%%% ACTIVITIES
\begin{minipage}[t]{0.1\linewidth}
\textbf{Activities}
\end{minipage}\hspace{0.05\linewidth}
\begin{minipage}[t]{0.8\linewidth}
\begin{mysec}
    2024 \>\> Berkeley Economists for Equity (BEE) Mentor\\ 
    2023 \>\> Department of Agricultural and Resource Economics Graduate Admissions Committee \\
    2022 \>\> Department of Agricultural and Resource Economics Search Committee \\
    2020 - 2022 \>\> Department of Agricultural and Resource Economics DEI committee \\


\end{mysec}
\end{minipage}\vspace{5mm}


%%% PRIOR EMPLOYMENT
\begin{minipage}[t]{0.1\linewidth}
\textbf{Prior \\ Employment}
\end{minipage}\hspace{0.05\linewidth}
\begin{minipage}[t]{0.8\linewidth}
\textbf{UC Berkeley}, Graduate Student Researcher (Marco Gonzalez-Navarro) \hfill 2023-2025 \\
\textbf{UC Berkeley}, Graduate Student Researcher (Ethan Ligon) \hfill 2022 \\
\textbf{UC Berkeley}, Graduate Student Researcher (Ellen Bruno) \hfill 2023 \\
\textbf{UC Berkeley}, Graduate Student Researcher (Joseph Shapiro, Andrés Rodríguez-Clare) \hfill 2021 \\
\textbf{National Laboratory for Public Policy (LNPP)}, Analyst \hfill 2018-2020 \\
\end{minipage}\vspace{5mm}

%%% ANY OTHER ITEMS
\begin{minipage}[t]{0.1\linewidth}
\textbf{}
\end{minipage}\hspace{0.05\linewidth}
\hspace{0.05\linewidth}
\begin{minipage}[t]{0.8\linewidth}


\end{minipage}\vspace{5mm}


\end{document}